\section{Conclusion}

The 2018 Hilbert-mapping paper should not be revived as though the field had stood still. Its empirical framing belongs to an earlier stage of 3D deep learning. Yet the paper's central instinct aged remarkably well: the representation used to expose structure to a model can matter as much as the model itself.

This upgraded manuscript argues that serialization is best understood as a first-class interface between irregular 3D structures and sequence-native learning systems. Under that view, the source paper becomes an early precursor to a modern family of ideas spanning serialized point transformers, state-space point and voxel models, neural surface reconstruction, and biomolecular structure learning. The resulting contribution is not a facelift. It is a distinct perspective paper with a new taxonomy, a new design framework, and a concrete publication agenda.

The broader lesson is simple. When a field rediscovers an old idea under new system constraints, the right response is not always to rerun the old benchmark. Sometimes the right response is to clarify the design principle the old paper noticed before the rest of the field had a name for it.
